\section{Panorama}

AgroVoltaic es la plataforma de datos del sistema agrovoltaico del TEC en San Carlos
(Costa Rica): paneles solares bifaciales instalados sobre cultivo. El repositorio es un
\emph{monorepo} con cuatro piezas que se pueden levantar por separado.

\begin{center}
\small
\begin{tabular}{@{}p{34mm}p{34mm}p{72mm}@{}}
\toprule
\textbf{Pieza} & \textbf{Dónde vive} & \textbf{Qué resuelve} \\
\midrule
ETL de CSV & \cd{src/agrovoltaic} & Estandarizar 19 meses de CSV crudos del inversor
y los sensores, y cargarlos a Supabase de forma idempotente. \\
Agente de pronóstico & \cd{agente-pronostico/} & Pronosticar irradiancia y humedad de
suelo con anclaje físico; el LLM sólo orquesta. \\
Agente analizador & \cd{agente-analizador/} & Responder preguntas sobre el histórico
fotovoltaico ya limpio, componiendo herramientas atómicas. \\
Consola de depuración & \cd{mvp-debugger/} & Ver la traza completa de cada agente y
cruzar cada número contra los datos. \\
\bottomrule
\end{tabular}
\end{center}

\subsection{El problema de origen}

Los datos no tenían estándar. El EDA sobre los 285 CSV encontró \textbf{13 esquemas
distintos}: nombres de columna que cambian entre cortos (\cd{vpv1}), largos
(\cd{Voltaje PV1 [V]}) y \cd{snake\_case}; erratas persistentes (\cd{Energì} con acento
grave en 72 archivos, \cd{POTencia}, \cd{Corriente PV2[A]} sin espacio); filas de
sensores distintos intercaladas en el mismo archivo; irradiancia sin calibrar con un
\emph{offset} nocturno constante de $-38{,}845$; temperaturas saturadas en 85\,°C
(el código de error clásico del DS18B20 desconectado); intervalos de muestreo que van
de 2 segundos a 5 minutos según la época; y dos huecos largos de 126 y 71 días.

\begin{clave}[title=Regla rectora del tratamiento de datos]
Validada con Leo Cardinale el 2026-08-10: \textbf{el dato crudo se guarda tal cual en la
base; toda corrección vive en una capa de análisis} (vistas SQL) que genera variables
nuevas. Ninguna transformación destructiva en la ingesta. Esta regla reemplazó las
decisiones previas de convertir 85\,°C a NULL, forzar el \emph{offset} a cero y
remuestrear todo a 5 minutos durante la carga.
\end{clave}

\subsection{Dos regiones, sin base central}

El proyecto abarca dos sitios y \textbf{no} se fusionan sus bases:

\begin{itemize}[leftmargin=*,itemsep=2pt]
  \item \textbf{Cartago} --- recolección automática hacia \textbf{AgroDash} (PostgreSQL,
    app Rust/Axum). Sensores de suelo y ambiente: humedad, conductividad, temperatura,
    irradiancia, PAR. No mide fotovoltaica.
  \item \textbf{San Carlos} --- recolección semi-manual de CSV hacia \textbf{Supabase}.
    Es la región fotovoltaica.
\end{itemize}

El matiz que define la arquitectura: \textbf{San Carlos está partido entre las dos
bases}. Lo eléctrico (voltaje, corriente, potencia, energía, irradiancia del inversor)
vive en Supabase; lo ambiental y de suelo de San Carlos vive en AgroDash, en cajas con
sufijo \cd{SC}. Como Cartago no mide fotovoltaica, no existe comparación cruzada
PV contra PV: lo único común entre regiones son las variables ambientales.

% Pagina propia, sin rotar: el diagrama esta dibujado en vertical justamente para
% que entre derecho y se lea sin ladear la cabeza.
\begin{figure}[p]
\centering
\resizebox{\textwidth}{!}{% Figura: arquitectura global. Flujo de arriba hacia abajo, una banda por etapa:
% fuentes -> acarreo -> ingesta -> almacen -> agentes -> clientes. Vertical a
% proposito: asi entra en una pagina normal sin rotarla ni encogerla hasta
% volverla ilegible.
\begin{tikzpicture}[
  banda/.style={font=\scriptsize\bfseries, color=humo, anchor=east},
  fila/.style={text width=27mm, minimum height=11mm},
]

% --- Coordenadas: una y por banda, una x por columna -------------------------
\def\yfuentes{0}
\def\yacarreo{-3.2}
\def\yingesta{-6.4}
\def\yalmacen{-9.6}
\def\yagentes{-12.8}
\def\yclientes{-16.0}
\def\xpv{3.2}
\def\xamb{10.6}

% --- Banda 1: fuentes fisicas ------------------------------------------------
\node[fuente,fila] (inv) at (0,\yfuentes) {Inversor solar\\\nota{PV1 + PV2}};
\node[fuente,fila] (pir) at (\xpv,\yfuentes) {Piranómetros\\\nota{incidente / reflejada}};
\node[fuente,fila] (ds)  at (6.4,\yfuentes) {DS18B20\\\nota{temp.\ de arreglo}};
\node[agro,fila] (cajas) at (\xamb,\yfuentes) {Cajas y nodos SC\\\nota{suelo, ambiente}};

\node[grupo, fit=(inv)(pir)(ds),
      label={[titgrupo]above:SAN CARLOS --- FOTOVOLTAICO}] (gf) {};
\node[grupo, fit=(cajas),
      label={[titgrupo]above:SAN CARLOS --- AMBIENTAL}] (gc) {};

% --- Banda 2: acarreo --------------------------------------------------------
\node[comp,fila] (csv) at (\xpv,\yacarreo)
  {CSV crudos\\\nota{285 archivos}\\\nota{13 esquemas}};
\node[agro,fila] (agrodash) at (\xamb,\yacarreo)
  {AgroDash\\\nota{PostgreSQL}\\\nota{región Cartago}};

% --- Banda 3: ingesta --------------------------------------------------------
\node[comp,fila] (etlcsv) at (\xpv,\yingesta)
  {\textbf{ETL de CSV}\\\cd{src/agrovoltaic}\\\nota{lote, manual}};
\node[comp,fila] (etlamb) at (\xamb,\yingesta)
  {\textbf{ETL ambiental}\\\cd{pronostico.etl}\\\nota{timer, 15 min}};

% --- Banda 4: almacenamiento -------------------------------------------------
\node[almacen,text width=46mm] (store) at (6.9,\yalmacen)
  {\textbf{Supabase AgroVoltaic}\\\nota{crudo + vistas + store del agente}};

% --- Banda 5: agentes --------------------------------------------------------
\node[comp,fila] (ana) at (\xpv,\yagentes)
  {\textbf{Agente analizador}\\\nota{Q\&A histórico PV}\\\nota{\cd{:8010}}};
\node[comp,fila] (pro) at (\xamb,\yagentes)
  {\textbf{Agente de pronóstico}\\\nota{irradiancia + humedad}\\\nota{\cd{:8000}}};
\node[compx,text width=22mm] (llm) at (6.9,\yagentes) {API de Claude};

% --- Banda 6: clientes -------------------------------------------------------
\node[comp,fila] (web) at (4.6,\yclientes)
  {\textbf{Consola}\\\cd{mvp-debugger}\\\nota{Next.js, local}};
\node[compx,fila] (vf) at (9.2,\yclientes) {VisioneFlow\\\nota{flujos programados}};

% --- Etiquetas de banda ------------------------------------------------------
\node[banda] at (-1.9,\yfuentes)  {FUENTES};
\node[banda] at (-1.9,\yacarreo)  {ACARREO};
\node[banda] at (-1.9,\yingesta)  {INGESTA};
\node[banda] at (-1.9,\yalmacen)  {ALMACÉN};
\node[banda] at (-1.9,\yagentes)  {AGENTES};
\node[banda] at (-1.9,\yclientes) {CLIENTES};

% --- Flujos ------------------------------------------------------------------
\draw[flecha] (inv.south) -- (csv.north);
\draw[flecha] (pir.south) -- (csv.north);
\draw[flecha] (ds.south)  -- (csv.north);
\draw[flecha] (cajas)     -- (agrodash);

\draw[flecha] (csv)      -- (etlcsv);
\draw[flecha] (agrodash) -- node[etiq,right]{sólo lectura} (etlamb);

\draw[flecha] (etlcsv.south) -- ++(0,-0.6) -| (store.north);
\draw[flecha] (etlamb.south) -- ++(0,-0.6) -| (store.north);

\draw[flecha] (store.south) -- ++(0,-0.6) -| (ana.north);
\draw[flecha] (store.south) -- ++(0,-0.6) -| (pro.north);
\node[etiq,anchor=south west] at ($(ana.north)+(0.12,0.02)$) {sólo lectura};
\node[etiq,anchor=south east] at ($(pro.north)+(-0.12,0.02)$) {lee y escribe predicciones};

\draw[flechad] (ana) -- (llm);
\draw[flechad] (pro) -- (llm);

\draw[flecha]  (ana.south) -- ++(0,-0.6) -| (web.north);
\draw[flecha]  (pro.south) -- ++(0,-0.6) -| (web.north);
\draw[flechad] (pro.south) -- ++(0,-0.6) -| (vf.north);

\end{tikzpicture}
}
\caption{Arquitectura global. Dos rutas de ingesta independientes convergen en una sola
Supabase, que es lo único que los agentes leen.}
\end{figure}
